\section{Results}\label{sec:results}

\subsection{Sample and reconstruction (RQ1)}\label{sec:res_recon}

The primary sample contains \nPlayers{} players and \nObs{} valid PER observations (Table~\ref{tab:flow}); the median career has \medianSeasons{} valid seasons (mean \meanSeasons), the longest 23, and positions are balanced (318 to 349 players per modal position). Figure~\ref{fig:examples} shows four careers that illustrate the heterogeneity the analysis must accommodate, including gaps.

\begin{table}[t]\centering\small
\caption{Construction of the analysis sample under the primary and alternative specifications. Players enter if their first NBA season is in the debut window; ``career games'' and ``valid seasons'' are the successive restrictions; the last column counts season observations that enter the loss.}
\label{tab:flow}
\resizebox{\textwidth}{!}{\begin{tabular}{lrrrr}\toprule
Specification & Debut in window & Career games & Valid seasons & PER observations \\ \midrule
Primary & 2,381 & 1,833 & 1,676 & 13,193 \\
No minutes threshold, PER$>$45 removed (2025 version) & 2,381 & 1,833 & 1,749 & 14,128 \\
250-minute threshold & 2,381 & 1,833 & 1,544 & 12,168 \\
At least 3 valid seasons & 2,381 & 1,833 & 1,444 & 12,729 \\
No career-games restriction & 2,381 & 2,381 & 1,685 & 13,211 \\
Debut 1979--80 to 2005--06 & 2,051 & 1,577 & 1,437 & 11,254 \\
\bottomrule\end{tabular}}
\end{table}

\begin{figure}[t]\centering
\includegraphics[width=\textwidth]{fig1_examples}
\caption{Four PER careers on the calendar axis. Shaded seasons are gaps (fewer than 100 minutes); the dotted line is the league average of 15. Michael Jordan's two retirements produce gaps that are masked rather than interpolated.}
\label{fig:examples}
\end{figure}

The autoencoder with $d=64$ reaches a held-out RMSE of \rmseHeld{} PER points (SD over five seeds \rmseHeldSd; training RMSE \rmseTrain), against \rmseLinear{} for a player-specific linear trend, \rmsePlayerMean{} for the player's own mean and \rmseGlobal{} for the global mean (Table~\ref{tab:recon}, Figure~\ref{fig:recon}); the prespecified criterion is met by a wide margin, and a 70/10/20 split gives \rmseTestThreeWay{} (Appendix~\ref{app:training}). Error falls steeply from $d=8$ (\rmseHeldEight) to $d=32$ (\rmseHeldThirtytwo) and is flat thereafter (\rmseHeldOneTwoEight{} at $d=128$). Because no career exceeds \maxSpan{} seasons, a code of 32 or more dimensions can store the sequence itself; the plateau shows that the model does so almost losslessly. This matters for clustering: a near-lossless code preserves all the idiosyncratic detail of a career, including detail shared with no other career.

\begin{table}[t]\centering\small
\caption{Reconstruction RMSE in PER units on the development (train) and held-out players, by latent dimensionality, with three baselines evaluated on the same held-out observations.}
\label{tab:recon}
\resizebox{\textwidth}{!}{\begin{tabular}{lcccc}\toprule
Model & Seeds & Train RMSE & Held-out RMSE (SD over seeds) & Best epoch \\ \midrule
Transformer autoencoder, $d=8$ & 3 & 0.88 & 1.23 (0.07) & 309 \\
Transformer autoencoder, $d=16$ & 3 & 0.61 & 0.95 (0.20) & 353 \\
Transformer autoencoder, $d=32$ & 3 & 0.47 & 0.80 (0.08) & 367 \\
Transformer autoencoder, $d=64$ & 10 & 0.43 & 0.74 (0.06) & 361 \\
Transformer autoencoder, $d=128$ & 3 & 0.41 & 0.74 (0.02) & 358 \\
\midrule Player-specific linear trend (2 parameters) & -- & -- & 2.24 & -- \\
Player mean & -- & -- & 2.76 & -- \\
Global mean & -- & -- & 4.35 & -- \\
\bottomrule\end{tabular}}
\end{table}

\begin{figure}[t]\centering
\includegraphics[width=0.55\textwidth]{fig3_reconstruction}
\caption{Held-out reconstruction RMSE of the Transformer autoencoder as a function of latent dimensionality (mean and SD over seeds), with the three baselines.}
\label{fig:recon}
\end{figure}

\subsection{A single UMAP+HDBSCAN run is not reproducible (RQ2)}\label{sec:res_single}

Over the \nGrid{} grid configurations, HDBSCAN returned between \gridKmin{} and \gridKmax{} clusters. The three selection rules chose three different solutions (Appendix Table~\ref{tab:selection}): the DBCV rule a ten-dimensional embedding with \selDbcvKdev{} clusters on the development set, \selDbcvKho{} on the held-out players and \selDbcvKall{} when refitted on all players; the Silhouette rule and the composite rule of the earlier version two-dimensional embeddings with \selSilKdev{} and \selCompKdev{} clusters, which became \selSilKall{} and \selCompKall{} on refitting, with \selSilNoiseho\% and \selCompNoiseho\% of held-out players unassigned. The selected solution depends more on the rule and the sample than on the data.

Figure~\ref{fig:single} (details in Appendix Table~\ref{tab:stability}) measures reproducibility directly. Holding the autoencoder fixed and changing only the UMAP random state, the mean pairwise ARI between runs ranged from \ariSingleMin{} to \ariSingleMax{} across configurations and extraction methods, never reaching the criterion of 0.80, and the number of clusters within a configuration varied between \kSingleMin{} and \kSingleMax{}. With excess-of-mass extraction the largest cluster contained between \bigShareEomMin\% and \bigShareEomMax\% of players on average, and it was this cluster whose subdivision changed from run to run; leaf extraction produced more and smaller clusters at the cost of \noiseLeafMean\% noise. Changing the autoencoder seed at a fixed UMAP seed gave mean ARIs between \ariAeSeedsMin{} and \ariAeSeedsMax{}. RQ2 is answered in the negative: the design of the earlier version, and by extension its nineteen-cluster solution, does not describe a reproducible partition. The instability has two sources, the stochastic embedding and, more fundamentally, the dependence of the learned latent geometry on the training seed.

\begin{figure}[t]\centering
\includegraphics[width=0.9\textwidth]{fig4_single_run}
\caption{Reproducibility of single UMAP+HDBSCAN runs across ten UMAP random states (three autoencoder seeds each) for four configurations (C1--C4, Appendix Table~\ref{tab:stability}) and two extraction methods. Left: pairwise ARI (mean, with range); the dotted line is the prespecified criterion. Right: number of clusters (median, with range).}
\label{fig:single}
\end{figure}
